\section{Some Motivation}

\subsection{Degrees of Maps on $n$-Spheres}

\begin{boxdefinition}[Degree]
    Let $f$ be a continuous map from $S^n$ to itself. We know there is an induced map $f_{*} : H_n\of{S^n} \to H_n\of{S^n}$. We know that $H_n\of{S^n} \cong \Z$. Factoring through this isomorphism, we call $f_{*}(1)$ the \textbf{degree} of $f$, denoted $\pdeg{f}$.
\end{boxdefinition}

The intuition behind this number is that it tells you ``how many times your map $f$ wraps the $n$-sphere around itself.''

\begin{boxlemma}[Facts About the Degree] \hfill
    \begin{enumerate}
        \item If $f$ is not surjective then $\pdeg{f} = 0$.
        \item $\pdeg{f \circ g} = \pdeg{f} \cdot \pdeg{g}$.
        \item $\pdeg{\id} = 1$.
        \item If $f \simeq g$ (ie, $f$ is homotopic to $g$), then $\pdeg{f} = \pdeg{g}$.\footnote{The converse is \textit{also} true: it's called Hopf's Theorem.}
        \item If $f$ is a reflection of some map via a hyperplane, then $\pdeg{f} = -1$.
    \end{enumerate}
\end{boxlemma}
\begin{proof}[Proof Sketch] \hfill
    \begin{enumerate}
        \item If $f$ is not surjective, then there is some $x_0 \in S^n$ such that $f$ corestricted to $S^n \setminus \set{x_0}$ agrees with $f$. That is, the following diagram commutes:
        \begin{cd*}
            S^n \arrow[rr, "f"] \arrow[dr, "f"'] && S^n \\[1em]
            & S^n \setminus \set{x_0} \cong B^n \arrow[ur, hook]
        \end{cd*}
        By functoriality of homology, we know that taking $H_n$ preserves the above commutative diagram. \sorry

        \item This follows from the functoriality of homology.

        \item This also follows from the functoriality of homology.

        \item Something to do with agreement of homology on homotopic spaces/maps? \sorry

        \item \sorry
    \end{enumerate}
\end{proof}

To see why this is useful, we consider groups acting on $n$-spheres.

Recall the notion of a \textbf{free action}.

\begin{boxdefinition}[Free Group Action]
    A group $G$ \textbf{acts freely} on a set $X$ if the action has no fixed points, ie, if every nontrivial element of $G$ fixes no elements of $X$.
\end{boxdefinition}

It turns out maps with no fixed points have a special property.

\begin{boxlemma}
    If $f : S^n \to S^n$ is continuous and has no fixed points, then $f$ is homotopy equivalent to the antipodal map $a : S^n \to S^n : x \mapsto -x$. In particular, $\pdeg{f} = \parenth{-1}^{n+1}$.
\end{boxlemma}
\begin{proof}[Proof Sketch]
    We can define an explicit homotopy $F : S^n \times I \to S^n$ by
    \begin{align*}
        F(x, t)
        = \frac{
            t \cdot f(x) + \parenth{1-t} \cdot \parenth{-x}
        }{
            \norm{t \cdot f(x) + \parenth{1-t} \cdot \parenth{-x}}
        }
    \end{align*}
    It is not difficult to show that this is a well-defined homotopy.
\end{proof}

We can now use the machinery about degrees to show the following rather nontrivial fact.

\begin{boxtheorem}
    Let $n$ be an even number. The only nontrivial group that acts freely on $S^n$ is $\Zmod{2}$.
\end{boxtheorem}
\begin{proof}
    Let $G$ be a group acting freely on $S^n$. Then,
    \begin{align*}
        \deg : G \to \set{\pm 1} \cong \Zmod{2}
    \end{align*}
    is a group homomorphism.
    
    Since the action is free, every $g \in G \setminus \set{\id}$ is a homeomorphism without fixed points. \sorry % (Apply Lemma in some way)
\end{proof}

Note that the above is really not true of odd-dimensional spheres.

\begin{boxcexample}[Many Free Actions on $S^3$]
    We can view $S^3$ as a subset of $\C^2$. We can show that $\Zmod{n}$, generated by some $\epsilon$, acts freely on $S^3$. Indeed, map any pair $\parenth{z, w} \in \C^2$ to the point
    \begin{align*}
        \parenth{e^{\frac{2 \pi i}{n}} \cdot z, e^{\frac{2 \pi i}{n}} \cdot w}
    \end{align*}
    In other words, we rotate both coordinates in some special way. This action has no fixed points when restricted to $S^3$. Fun fact: the quotient of $S^3$ by the above action is called a \textbf{lens space}.
\end{boxcexample}

We can do something similar for any odd-dimensional sphere, as it lives inside a Euclidean space of even dimension.

\begin{boxtheorem}
    $S^n$ has a nonzero tangent vector field (that is, $\setst{v : S^n \to \R^{n + 1}}{\cycl{v(x), x} = 0} \neq 0$) if and only if $n$ is odd.
\end{boxtheorem}
\begin{proof}
    \sorry
\end{proof}

Finally, we outline some useful homology theoretic machinery for computing degrees.

It turns out the following is a short exact sequence of chain complexes:
\begin{cd*}
    0 \arrow[r] &
    C_{\bullet}\of{\R P^n; \Zmod{2}} \arrow[r, "\tau"] &
    C_{\bullet}\of{S^n; \Zmod{2}} \arrow[r, "q_{\bullet}"] &
    C_{\bullet}\of{\R P^n; \Zmod{2}} \arrow[r] &
    0
\end{cd*}
where $q : S^n \to \R^n$ is given by mapping $\sigma$ to $\tilde{\sigma_1} + \tilde{\sigma_2}$, where $\tilde{\sigma_1}$ and $\tilde{\sigma_2}$ are the two lifts of $\sigma$.